\subsection{Poincaré group of a graph}

Let G be a graph. Let a be a vertex of G; the isotropy group at a of the groupoid \(\mathrm{Grp}(G)\) is called the Poincaré group of G at a and is denoted \(\pi_{1}(G,a)\) . Let c be a path class in G, let a be its origin and b its target. The map \(\operatorname{Int}(c): \pi_{1}(G,b)\to\pi_{1}(G,a)\) defined by \(c'\mapsto cc'c^{-1}\) is a group isomorphism. Let \(\varphi:G\to H\) be a graph morphism. Denote by \(\theta_{G}:G\to\mathrm{Grp}(G)\) and \(\theta_{H}:H\to\mathrm{Grp}(H)\) the canonical morphisms. If \(\overline{\varphi}\) denotes the unique groupoid morphism \(\mathrm{Grp}(G) \to \mathrm{Grp}(H)\) such that \(\overline{\varphi}\circ\theta_{G}=\theta_{H}\circ\varphi\) , the group homomorphism \(\overline{\varphi}_{a}:\pi_{1}(G,a)\to\pi_{1}(H,\varphi(a))\) is denoted \(\pi_{1}(\varphi,a)\) .

Let G be a connected graph, let S be an orientation of G, and let A be a maximal oriented tree of G (II, p. 157, prop. 1). Given vertices \(a\) and \(b\) of G, there exists a unique path class \(\gamma_{a,b}\) in the graph \(\widetilde{\mathrm{A}}\) associated with A which connects \(a\) to \(b\) (II, p. 174, cor. 4 of prop. 5). If \(a\) , \(b\) and \(c\) are vertices of G, we have \(\gamma_{a,b}\gamma_{b,c} = \gamma_{a,c}\) , these two classes of paths being equal to the unique class of paths connecting \(a\) to \(c\) in \(\widetilde{\mathrm{A}}\) .

PROPOSITION 9. --- Let a be a vertex of G. There exists a unique homomorphism \(\lambda\) from the free group \(\mathrm{F}(\mathrm{S}-\mathrm{Fl}(\mathrm{A}))\) to \(\pi_1(G, a)\) such that

\[
\lambda (f) = \gamma_ {a, o (f)} \cdot f \cdot \gamma_ {t (f), a}\tag{1}
\]

for \(f \in \mathrm{S} - \mathrm{Fl}(\mathrm{A})\) . The homomorphism \(\lambda\) is a group isomorphism.

Let L denote the group \(\mathrm{F}(\mathrm{S}-\mathrm{Fl}(\mathrm{A}))\) . The existence and uniqueness of the homomorphism \(\lambda\colon L\to\pi_{1}(\mathrm{G},a)\) satisfying relations (1) follows from A, I, p. 85, prop. 8. Let L be the groupoid associated with the group L (II, p. 163, example 1); denote by s its unique vertex. There exists a unique groupoid morphism \(\mu\colon\mathrm{Grp}(\mathrm{G})\to\mathcal{L}\) such that \(\mu(\theta_{\mathrm{G}}(f))=f\) for every arrow \(f\in\mathrm{S}-\mathrm{Fl}(\mathrm{A})\) and \(\mu(\theta_{\mathrm{G}}(f))=e_{s}\) for every arrow \(f\in\mathrm{Fl}(\mathrm{A})\) (II, p. 173, prop. 5). For \(f\in\mathrm{S}-\mathrm{Fl}(\mathrm{A})\) , we have \(\mu(\theta_{\mathrm{G}}(f))=f\) ; the homomorphism \(\mu_{a}\circ\lambda\) is therefore the identity isomorphism of the group L (A, I, p. 85, prop. 8). It follows that \(\lambda\) is injective.

Let \(\varpi_{\mathrm{G}} / \mathrm{A}\) be the groupoid deduced from \(\varpi_{\mathrm{G}}\) by contraction of the arrows of A and let \(p\colon \varpi_{\mathrm{G}}\to \varpi_{\mathrm{G}} / \mathrm{A}\) be the canonical groupoid morphism. It is surjective, and the group homomorphism \(p_a\) deduced from \(p\) by passing to isotropy groups is an isomorphism (II, p. 178, cor. 2 of prop. 8). Since the graph G is assumed connected and A is a maximal tree thereof, the groupoid \(\varpi_{\mathrm{G}} / \mathrm{A}\) has a unique vertex (II, p. 176). Moreover, \(\varpi_{\mathrm{G}}\) is generated by \(\theta_{\mathrm{G}}(\mathrm{G})\) ; therefore, \(\varpi_{\mathrm{G}} / \mathrm{A}\) is generated by the loops \(p(f)\) , for \(f\in \mathrm{S}\) , and the group \((\varpi_{\mathrm{G}} / \mathrm{A})_{p(a)}\) is generated by the elements of the form \(p(\theta_{\mathrm{G}}(f))\) , for \(f\in \mathrm{S} - \mathrm{Fl}(\mathrm{A})\) . The homomorphism \(p_a\circ \lambda\) is therefore surjective. Consequently, \(\lambda\) is surjective.

Remarks. --- 1) The homomorphism \(\mu_{a}\) is the inverse isomorphism of \(\lambda\) .

\begin{enumerate}
\def\labelenumi{\arabic{enumi})}
\setcounter{enumi}{1}
\tightlist
\item
  There exists a unique homomorphism \(\lambda\colon\mathrm{F}(\mathrm{Fl}(\mathrm{G}))\to\pi_{1}(\mathrm{G},a)\) defined by relations (1) for all \(f\in\mathrm{Fl}(\mathrm{G})\) . It follows from Proposition 8 that the homomorphism \(\lambda\) is surjective and its kernel is the smallest normal subgroup of \(\mathrm{F}(\mathrm{Fl}(\mathrm{G}))\) containing the elements \(f\) , for \(f\in\mathrm{Fl}(\mathrm{A})\) , and the elements \(f\cdot\overline{f}\) , for \(f\in\mathrm{Fl}(\mathrm{G})\) .
\end{enumerate}

